\documentclass[11pt,a4paper]{article}
\usepackage[margin=25mm]{geometry}
\usepackage{graphicx}
\usepackage{amsmath,amssymb}
\usepackage{booktabs,longtable,array}
\usepackage{hyperref}
\usepackage{enumitem}
\usepackage{pdflscape}
\usepackage{tikz}
\usetikzlibrary{arrows.meta}
\setlist{nosep}
\hypersetup{colorlinks=true,linkcolor=blue,urlcolor=blue}
\title{Spatially Correlated Ice-Cliff Presence and Area-Fraction Prediction in IGM}
\author{Method and IGM User-Process Implementation}
\date{20 October 2026}

\begin{document}
\maketitle
\tableofcontents
\bigskip

\clearpage
\begin{landscape}
\section{End-to-End Pipeline}
\begin{center}
\begin{tikzpicture}[
  trainbox/.style={draw=blue!65!black,fill=blue!5,rounded corners=2pt,
    text width=4.25cm,minimum height=3.15cm,inner sep=5pt,align=left,
    font=\scriptsize},
  runbox/.style={draw=green!45!black,fill=green!5,rounded corners=2pt,
    text width=4.25cm,minimum height=3.15cm,inner sep=5pt,align=left,
    font=\scriptsize},
  outputbox/.style={draw=orange!65!black,fill=orange!7,rounded corners=2pt,
    text width=6.2cm,minimum height=3.15cm,inner sep=6pt,align=left,
    font=\scriptsize},
  flow/.style={-{Stealth[length=2mm]},draw=black!70,line width=.7pt},
  artifact/.style={-{Stealth[length=2mm]},draw=blue!65!black,
    dashed,line width=.7pt},
  phase/.style={font=\small\bfseries,anchor=west}
]
\node[phase,text=blue!65!black] at (.25,13.1)
  {OFFLINE TRAINING AND PARAMETER FITTING};
\node[trainbox] (train1) at (2.5,11.2) {
  \textbf{Training targets and split}\\
  At 100 m, a category is present when its fraction exceeds $10^{-8}$. Sample
  debris-positive cells; train on 50 glaciers and hold out 17. Held-out
  glaciers do not contribute to fitting, calibration, or covariance estimates.
};
\node[trainbox] (train2) at (7.45,11.2) {
  \textbf{Feature construction}\\
  Four predictors: velocity, driving stress, debris thickness, and slope.
  Build local center and 3$\times$3 / 11$\times$11 statistics, plus upglacier
  mean, minimum, maximum, and trend at 200, 500, and 1000 m. The full vector
  has 68 features; crevasses use 20 window features.
};
\node[trainbox] (train3) at (12.35,11.2) {
  \textbf{Category classifiers}\\
  Fit one binary random forest per category (run 015): 500 trees, minimum leaf
  size 5, maximum feature fraction 0.5, balanced bootstrap samples, and median
  imputation. Upglacier traces use observed velocity and fall back to surface
  descent where needed.
};
\node[trainbox] (train4) at (17.25,11.2) {
  \textbf{Probability and presence covariance}\\
  Generate five-fold glacier-grouped out-of-fold predictions. Fit a separate
  monotone Platt calibration per category. Standardized residuals estimate
  category-specific squared-exponential presence ranges.
};
\node[trainbox] (train5) at (22.15,11.2) {
  \textbf{Fraction distributions and artifacts}\\
  Fit each positive-fraction marginal by discretized likelihood and AIC. Fit
  magnitude covariance from within-glacier normal scores. Save the forests in
 \path{cliff_prediction_models.joblib} and calibration, ranges, and
 distributions in \path{cliff_prediction_parameters.json}.
};
\draw[flow] (train1.east) -- (train2.west);
\draw[flow] (train2.east) -- (train3.west);
\draw[flow] (train3.east) -- (train4.west);
\draw[flow] (train4.east) -- (train5.west);

\node[phase,text=green!45!black] at (.25,8.15)
  {RUNTIME INFERENCE: LOAD THE RUN-015 MODEL AND PARAMETER BUNDLE};
\node[runbox] (run1) at (2.5,6.2) {
  \textbf{IGM state and predictors}\\
  Read debris thickness, slope, elevation, mask, and modeled
  \texttt{uvelsurf}/\texttt{vvelsurf}. Use modeled speed and direction; driving
  stress defaults to \texttt{dStress}, then \texttt{tau\_d}, then
  $\rho gH\sin(\alpha)$.
};
\node[runbox] (run2) at (7.45,6.2) {
  \textbf{Resample and build features}\\
  Convert predictor grids, velocities, elevation, and mask to 100 m. Compute
  local and upglacier features. Trace backward along smoothed modeled flow;
  use negative surface gradient when speed is low or missing.
};
\node[runbox] (run3) at (12.35,6.2) {
  \textbf{Predict category probabilities}\\
  Apply the four run-015 forests with training medians for missing features.
  Retain raw probabilities and, in stochastic mode, apply the saved Platt
  calibration independently for each category.
};
\node[runbox] (run4) at (17.25,6.2) {
  \textbf{Sample spatial presence}\\
  For each category, filter white noise to its fitted range, map through the
  normal CDF, then compare with calibrated probability to draw presence.
  Independent stochastic and raw-probability threshold modes are alternatives.
  Seeds derive reproducibly from glacier and category identifiers.
};
\node[runbox] (run5) at (22.15,6.2) {
  \textbf{Sample category fractions}\\
  Independent magnitude fields are mapped through each fitted inverse CDF;
  absent cells receive zero. Categories may overlap. If their sum exceeds one,
  the default rescales all fractions; alternatives allow the sum or raise an
  error.
};
\draw[flow] (run1.east) -- (run2.west);
\draw[flow] (run2.east) -- (run3.west);
\draw[flow] (run3.east) -- (run4.west);
\draw[flow] (run4.east) -- (run5.west);
\draw[artifact] (train5.south) -- ++(0,-.25) -|
  node[pos=.25,above,font=\tiny,fill=white,inner sep=1pt]
    {saved forests and parameters} (run3.north);

\node[phase,text=orange!65!black] at (.25,3.05)
  {NATIVE-GRID OUTPUT AND DEBRIS-SMB COUPLING};
\node[outputbox] (out1) at (4,1) {
  \textbf{Return to the IGM grid}\\
  Upsampling: probabilities use bilinear interpolation, fractions and presence
  use nearest-neighbor. Downsampling: use subcell area sampling; a native cell
  is present if any contributing 100 m cell is present.
};
\node[outputbox] (out2) at (12.35,1) {
  \textbf{Write state fields}\\
  Store stream, pond, crevasse, and undefined fractions in
  \texttt{state.cliffs\_s}, \texttt{cliffs\_p}, \texttt{cliffs\_c}, and
  \texttt{cliffs\_u}; store their sum in \texttt{state.cliffs\_total}.
};
\node[outputbox] (out3) at (20.7,1) {
  \textbf{Adjust debris SMB}\\
  When cliff SMB is enabled, blend debris and clean-ice SMB using total cliff
  fraction and \texttt{cliff\_smb\_enhancement}:\\
  $\mathrm{SMB}_{\rm out}=\mathrm{SMB}_{\rm debris}(1-f_{\rm cliff})+$\\
  $\mathrm{SMB}_{\rm clean}f_{\rm cliff}E_{\rm cliff}$.
};
\draw[flow] (out1.east) -- (out2.west);
\draw[flow] (out2.east) -- (out3.west);
\draw[flow] (run5.south) -- ++(0,-.35) -|
  node[pos=.25,above,font=\tiny,fill=white,inner sep=1pt]
    {native-resolution conversion} (out1.north);

\node[anchor=west,font=\tiny,text=black!70] at (.25,-1.05)
  {Default configuration: run 015, spatially stochastic presence, and fraction rescaling.
  The legacy three-predictor model is selected by its bundle filename.};
\end{tikzpicture}
\end{center}
\end{landscape}
\clearpage

\section{Purpose and Scope}
This document describes the cliff predictor installed in the IGM 3.2 user code at
\texttt{user/code/processes/cliff\_prediction.py}. It predicts four separate
cliff categories on a debris-covered glacier grid:
\begin{enumerate}
  \item stream-influenced,
  \item pond-influenced,
  \item crevasse-originated, and
  \item undefined.
\end{enumerate}
For each category, the method first predicts a presence probability, generates a
spatially correlated binary presence field, and then samples a positive cell-area
fraction for present cells. Category fractions are combined into a total cliff
area fraction. The public array API can be called by other IGM modules; the
\texttt{predict\_from\_state} adapter and the debris-SMB hook provide the current
IGM integration.

The deployed default classifier is run 015: the 100 m, four-predictor, per-category
presence random forests. The four predictors are \texttt{velocity\_Millan2022},
\texttt{dStress}, \texttt{debris\_thickness}, and \texttt{slope\_rad}. The model
and its parameters are stored in
\texttt{user/code/processes/cliff\_prediction\_models.joblib} and
\texttt{user/code/processes/cliff\_prediction\_parameters.json}.

Two additional classifier backends are available. Run 039 uses category-specific
50 m windows selected in run 037 and adds \texttt{elevation\_AW3D} only to the
pond model. Run 041 is the 100 m analogue, with upstream features for stream,
pond, and undefined; crevasse remains window-only. Both bundles are installed
in \texttt{user/code/processes/}. They reuse the default positive-fraction
marginals and spatial covariance parameters, but have no backend-specific Platt
calibrators, so probability calibration defaults off and must remain off unless
matching calibrators are fitted.

\section{Training Data and Classifier}
\subsection{Training target and domain}
At training time, a category is present in a 100 m cell when its raw category
area fraction is greater than $10^{-8}$. Training cells are debris-positive
cells with finite category targets. Half of each eligible glacier's cells were
sampled using the run-015 random seed. The glacier split held out 17 glaciers
and trained on 50 glaciers. The held-out glaciers were not used for classifier
fitting, probability calibration, or covariance estimation.

One binary random-forest classifier was fitted for each category. The run-015
parameters were 500 trees, minimum leaf size 5, maximum feature fraction 0.5,
class-balanced bootstrap samples, and random seed 42. Missing features are
imputed with training-set feature medians. The crevasse model uses only the
window features; the other three category models also use upstream features.

\subsection{Predictor features}
The training grid is 100 m. For each of the four predictors, the local window
features comprise:
\begin{itemize}
  \item the center value, 3\,$\times$\,3 finite-neighbour mean, and 3\,$\times$\,3 finite-neighbour standard deviation (12 values total);
  \item 11\,$\times$\,11 finite-neighbour mean and standard deviation, excluding the center (8 more values).
\end{itemize}
The remaining 48 features summarize values along an upglacier trace: four
statistics (mean, minimum, maximum, and trend) for each predictor at 200, 500,
and 1000 m. The trend is the least-squares slope with distance, expressed per
100 m, and includes the center value. Thus the full vector has 68 features;
the crevasse classifier selects 20 of them.

\subsection{Upglacier feature traces}
When the training features were prepared, Millan velocity components defined
the downstream direction. Traces move backward, upglacier, using a 50 m
midpoint integration step. Velocity components are Gaussian-smoothed over
about 100 m. Where speed is below 1 m\,yr$^{-1}$ or velocity is missing,
direction falls back to the negative gradient of a surface-elevation field
smoothed over about 150 m. Samples use bilinear interpolation and stop
contributing when direction or predictor values are unavailable. In the IGM
implementation, the direction is derived from modeled \texttt{uvelsurf} and
\texttt{vvelsurf}, as specified for this application.

\subsection{Training probabilities}
Let $p_{ic}^{\rm raw}$ be the random forest probability for cell $i$ and
category $c$. This is the forest's estimated probability that the category is
present, before calibration. The goal of calibration is to make probabilities
more representative of observed frequencies. For example, among cells assigned
a probability near 0.7, roughly 70\% should be present if those probabilities
are well calibrated. Calibration changes probability values, not the order of
cells from lower to higher predicted probability.

The calibration data are generated separately for each category using five-fold
\texttt{GroupKFold} cross-validation on the 50 training glaciers. In each fold,
all sampled cells from some glaciers are held out together. A forest with the
run-015 settings is refitted on the other folds and predicts probabilities for
the held-out glaciers. Repeating this process gives every sampled training cell
one out-of-fold (OOF) probability from a forest that did not train on that
glacier. Holding out whole glaciers avoids letting nearby or otherwise related
cells from the same glacier appear on both sides of a fold. These OOF
probabilities are used to fit the calibrator; they are not the probabilities
used directly at inference time. At inference, the saved calibrator is applied
to the raw probabilities from the final run-015 forest.

A separate monotone Platt transform is fitted for each category:
\[
 p_{ic}^{\rm cal} =
 \operatorname{logit}^{-1}\!\left(a_c+b_c\operatorname{logit}
 \left(\operatorname{clip}(p_{ic}^{\rm raw},10^{-5},1-10^{-5})\right)\right),
 \qquad b_c>0.
\]
Clipping keeps probabilities of exactly zero or one away from infinite
log-odds. The parameters are chosen by minimizing weighted binary log loss on
the OOF predictions. Each cell receives a weight inversely proportional to the
number of sampled cells from its glacier, so every glacier has the same total
weight even if some glaciers contribute many more cells. A small penalty on the
log-slope stabilizes the fit, and writing $b_c=\exp(\gamma_c)$ guarantees
$b_c>0$.

The intercept $a_c$ shifts the logits up or down. The slope $b_c$ controls how
strongly the raw logits are stretched: values below one pull probabilities
toward 0.5, while values above one push them farther from 0.5, before the
intercept shift. Because the slope is positive, the transform preserves the
ranking of cells. The fitted run-015 parameters are:
\begin{center}
\begin{tabular}{lrr}
\toprule
Category & $a_c$ (intercept) & $b_c$ (slope)\\
\midrule
Stream & -0.363762 & 0.789387\\
Pond & -0.668262 & 1.063949\\
Crevasse & -1.161831 & 0.656355\\
Undefined & -0.575136 & 0.717679\\
\bottomrule
\end{tabular}
\end{center}
These values are stored in the JSON parameter file. Calibrated probabilities
are used for stochastic presence sampling and for the residuals used to fit
spatial presence covariance. The deterministic threshold option instead uses
the raw random-forest probability at the configured threshold, 0.5 by default,
to preserve the original run-015 comparison.

\section{Spatially Correlated Presence}
The classifier gives each 100 m cell a probability of containing a cliff of a
given category. Drawing a separate yes/no result for every cell would treat
nearby cells as unrelated and could produce a speckled map. This section adds
spatial structure to those draws: nearby cells are more likely to have similar
outcomes, while each cell keeps the probability predicted for it. The
dependence is estimated separately for each cliff category.

\subsection{Training the presence covariance}
To estimate how strongly presence tends to cluster, we compare the observed
training labels with their calibrated probabilities. The calculation uses
out-of-fold predictions, so each training cell's probability comes from a
model that did not train on that cell's glacier. For category $c$ and cell $i$,
the standardized residual is
\[
 r_{ic}=\frac{y_{ic}-p_{ic}^{\rm cal}}
 {\sqrt{p_{ic}^{\rm cal}(1-p_{ic}^{\rm cal})}},
\]
where $y_{ic}$ is 1 when the category is present and 0 otherwise. A positive
residual means presence was observed more often than the probability predicted;
a negative residual means it was observed less often. Dividing by the
Bernoulli standard deviation makes errors at cells with different probabilities
more comparable. Residuals are clipped to $[-6,6]$ to limit the influence of
extreme values, then centered within each glacier to remove its average
prediction error. They are put back at their sampled grid locations so that
neighbouring cells can be compared.

For each category, correlations between residuals are grouped by their raster
offset (the row and column separation between cells), up to 12 cells, or
1.2 km. Offsets with fewer than 40 cell pairs are left out because their
correlation estimates would be based on too little data. A smooth,
direction-independent model is then fitted to describe how correlation changes
with distance:
\[
 \rho_c(d)=\exp\!\left[-\frac{1}{2}\left(\frac{d}{L_c^{\rm pres}}\right)^2\right].
\]
Here $d$ is the distance between cells and $L_c^{\rm pres}$ is the fitted
correlation range. It is a characteristic distance, not a boundary beyond
which correlation becomes zero: at $d=L_c^{\rm pres}$, the fitted correlation
is about 0.61, and it declines smoothly with distance. The isotropic model
assumes the same relationship in every direction.
The fitted run-015 presence ranges are:
\begin{center}
\begin{tabular}{lr}
\toprule
Category & $L_c^{\rm pres}$ (m)\\
\midrule
Stream & 78.4\\
Pond & 81.3\\
Crevasse & 152.2\\
Undefined & 70.4\\
\bottomrule
\end{tabular}
\end{center}
These ranges are around or below the 100 m grid spacing. They therefore
describe dependence on the discrete grid, rather than a sharply defined
physical patch size.

\subsection{Gaussian-copula sampling}
For each glacier and category, the sampler starts with random noise and applies
a spatial Gaussian filter. This produces a standard-normal field $Z_{ic}$ in
which nearby values tend to be similar. The standard-normal cumulative
distribution function $\Phi$ converts those values to a uniform field
$U_{ic}=\Phi(Z_{ic})$. The conversion keeps the spatial ordering of the field
but makes each cell's value uniformly distributed between zero and one. A
presence draw is made by comparing that value with the cell's calibrated
probability:
\[
 Y_{ic}=\mathbb{1}\{U_{ic}<p_{ic}^{\rm cal}\}.
\]
For example, a cell with probability 0.7 is present when its uniform value is
below 0.7. Because a uniform value falls below 0.7 in 70\% of repeated draws,
the cell's probability is preserved. Since nearby uniform values tend to be
similar, nearby cells are also more likely to land on the same side of their
probability thresholds. This makes presence spatially clustered without
changing the probability assigned to each cell. It does not force neighbours
to have the same outcome. Each category uses its own independent random field.
Random seeds are derived from a base seed, glacier identifier, and category
index, so the draws are reproducible.

The implementation uses a Gaussian filter on the discrete grid. Its smoothing
width is chosen so that the correlation produced on the grid matches the fitted
squared-exponential model. This calibration is important because the fitted
ranges are near or below one 100 m cell: a continuous-space approximation
alone would not describe the discrete filter accurately. The random field is
generated on a padded grid and then cropped to the glacier grid. Padding
prevents the filter from treating opposite edges of the array as neighbours.

\section{Spatially Correlated Category Fractions}
\subsection{Fitted marginal distributions}
Only positive 100 m category fractions from the 50 training glaciers are used
to fit these distributions. Candidate families are fitted with a discretized
likelihood that accounts approximately for the fraction-value lattice, and the
family with lowest AIC is selected independently for each category:
\begin{center}
\begin{tabular}{lll}
\toprule
Category & Selected family & Parameters\\
\midrule
Stream & Two-component truncated exponential & logit weight $-0.8816$; rates $103.2776,15.0357$\\
Pond & Two-component truncated exponential & logit weight $-0.7862$; rates $73.2004,16.5310$\\
Crevasse & Truncated Kumaraswamy & $a=0.8452,\ b=5.8086$\\
Undefined & Truncated generalized gamma & $a=5.8524,\ c=0.2899,\ \mathrm{scale}=4.6658\times10^{-5}$\\
\bottomrule
\end{tabular}
\end{center}
All fitted distributions are supported on $(0,1]$. For a two-component
truncated exponential, for example,
\[
 F(x)=w\frac{1-e^{-\lambda_1x}}{1-e^{-\lambda_1}}
 +(1-w)\frac{1-e^{-\lambda_2x}}{1-e^{-\lambda_2}},
 \qquad w=\operatorname{logit}^{-1}(\eta),\quad 0<x\leq1.
\]
\begin{figure}[htbp]
\centering
\includegraphics[width=\textwidth]{cliff_fraction_distributions_50m_100m.png}
\caption{Observed cumulative distributions of positive cliff fractions at 50 m and 100 m, with the best-AIC fitted distribution for each category and resolution. Dashed curves show fitted distributions; the dotted vertical line marks the $10^{-3}$ presence threshold. AIC is based on the discretized likelihood.}
\label{fig:cliff-fraction-distributions}
\end{figure}

\subsection{Training the magnitude covariance}
Presence and fraction magnitude are modeled with separate spatial fields. The
presence field decides which cells contain a category; the magnitude field
decides how large its positive fraction is. Magnitude covariance is therefore
estimated only from positive training fractions, separately for each category.
Cells with zero fraction are excluded because they are modeled by the presence
process, not by the positive-fraction distribution.

Raw fractions have strongly skewed, category-specific distributions.
Estimating spatial correlation directly from them would mix two effects: the
shape and scale of the fraction distribution, and the tendency for nearby cells
to have similar relative magnitudes. To separate these, the positive values are
ranked within each glacier and category. If a glacier has $n$ positive values
and a value has rank $k$, its rank is converted to a normal score
\[
 z_k=\Phi^{-1}\!\left(\frac{k-0.5}{n}\right),
\]
where $\Phi^{-1}$ is the standard-normal quantile function. Tied fractions
receive their average rank. The scores are then centered and scaled to have mean
zero and standard deviation one within that glacier. Ranking makes the score
pattern depend on which cells have relatively small or large fractions, rather
than on the category's particular units, skew, or range. Doing this separately
by glacier also prevents glaciers with systematically larger fractions from
setting the spatial relationship for the others.

The scores are placed back at their original 100 m raster cells; cells without
a positive target remain missing. For every raster offset up to 12 cells
(1.2 km), the method collects pairs of finite scores at that offset from all
training glaciers. Offsets are pooled into distance bins rounded to 0.1 cell,
including different directions; bins with fewer than 30 pairs are discarded.
This pooling treats spatial dependence as isotropic. The empirical correlation
for each distance is the normalized sum of products of the paired,
within-glacier standardized scores. Thus each qualifying cell pair contributes
to the estimate, and well-sampled distances are given more influence in the
fit.

The empirical correlations are fitted to the isotropic squared-exponential
model
\[
 \rho_c(d)=\exp\!\left[-\frac{1}{2}
 \left(\frac{d}{L_c^{\rm mag}}\right)^2\right],
\]
where $d$ is separation and $L_c^{\rm mag}$ is the category-specific magnitude
correlation range. The fit minimizes a weighted squared error between empirical
and modeled correlations; the weight is the square root of the number of pairs
at that distance. This single range is a compact approximation to the observed
correlogram, not a correlation cutoff: at $d=L_c^{\rm mag}$, the modeled
correlation is about 0.61, and it decays smoothly with distance. Because the
model is isotropic, it does not represent flow-aligned or other directional
structure. The fitted run-015 magnitude ranges are:
\begin{center}
\begin{tabular}{lr}
\toprule
Category & $L_c^{\rm mag}$ (m)\\
\midrule
Stream & 60.3\\
Pond & 45.6\\
Crevasse & 53.8\\
Undefined & 52.5\\
\bottomrule
\end{tabular}
\end{center}

These ranges describe dependence in the normal-score ranks of positive
fractions. They are not ranges fitted to raw fractions, and they should not be
read as the diameter of a typical cliff patch. All are below the 100 m cell
spacing, so their practical meaning is the strength of dependence on the
discrete grid, particularly between nearby cells; the exact relationship is
affected by the raster-lattice filter used during sampling.

\subsection{Sampling magnitudes}
Sampling combines the presence decision, the fitted magnitude covariance, and
the fitted positive-fraction distribution in sequence. First, the presence
process described above determines the binary indicator $Y_{ic}$ for each cell
and category. It uses the classifier's cell-specific probability, with spatial
dependence supplied by the separate presence field. The magnitude step does not
change this decision: it assigns a size only after presence has been
determined.

For each glacier and category, a new Gaussian field $G_{ic}$ is generated
independently of the presence field and of the other categories' magnitude
fields. Its spatial correlation uses the category-specific magnitude range
estimated above. When fraction spatial correlation is enabled, the discrete
Gaussian filter is calibrated to approximate that range on the 100 m grid; when
it is disabled, the cells receive independent Gaussian values. The field is
sampled at the debris-positive cell locations being predicted.

Applying the standard-normal CDF gives $V_{ic}=\Phi(G_{ic})$, uniform values
between zero and one. These values act as percentiles in the fitted
category-specific positive-fraction distribution $F_c$. Thus
\[
 f_{ic}=F_c^{-1}(V_{ic});
 \qquad f_{ic}=0\quad\text{when }Y_{ic}=0.
\]

The rank transformation used during training estimates only the spatial
dependence; it does not supply or resample the fraction values at runtime.
Instead, the inverse CDF turns each simulated percentile into a new fraction.
Because both transformations preserve order, nearby correlated Gaussian values
tend to produce nearby percentiles and hence similar fraction ranks, while the
inverse CDF restores the category's fitted, generally skewed fraction
distribution. Before the total-fraction constraint below, the positive draws
therefore follow $F_c$ marginally, while their spatial ranks are correlated.
Fraction magnitude is not forced to increase with presence probability because
the presence and magnitude fields are generated independently.

The four categories are sampled independently, so more than one can be present
in a cell. Their fractions are then checked jointly. If the sum exceeds one,
the default \texttt{rescale} policy divides each category fraction in that cell
by the sum, preserving which categories are present and making the total equal
to one. This necessary constraint changes the sampled magnitudes in overfull
cells, so their final marginal distributions are no longer exactly $F_c$. The
\texttt{allow} policy leaves overfull sums unchanged; \texttt{error} raises an
exception instead.

\section{IGM Resolution and State Integration}
The selected backend determines the model grid: run 039 uses 50 m, while run
015, run 016, and run 041 use 100 m. The IGM state adapter therefore:
\begin{enumerate}
  \item Resamples predictor grids, modeled velocity components, surface elevation, and mask to the selected backend's feature grid.
  \item Computes that backend's category-specific windows and upstream features, then runs the classifier and sampler at its model resolution.
  \item Resamples probabilities and category fractions back to the IGM grid. Probabilities use bilinear interpolation when upsampling. Fractions use nearest-neighbor when upsampling. Downsampling uses subcell area sampling for each output. Presence is nearest-neighbor when upsampling and set present when any contributing model-grid cell is present when downsampling.
\end{enumerate}
When \texttt{require\_100m\_grid} is enabled, direct array calls require the
selected backend's model resolution; the state adapter resamples automatically.
Its default predictor velocity is
$\sqrt{u_{\rm surf}^2+v_{\rm surf}^2}$ from modeled
\texttt{uvelsurf}/\texttt{vvelsurf}. Driving stress defaults to
\texttt{state.dStress}, then \texttt{state.tau\_d}, then
$\rho gH\sin(\alpha)$ with 3\,$\times$\,3 averaging. Upstream direction uses the
same modeled velocity components; low or missing speed uses the negative
surface gradient. Grid-axis signs convert map velocity into array row/column
directions.

Runs 039 and 041 were trained with observed \texttt{elevation\_AW3D} for the
pond category. At inference, the state adapter uses
\texttt{state.elevation\_AW3D} if present, otherwise modeled
\texttt{state.usurf} as a proxy; callers of the direct array API can supply
\texttt{predictors[elevation\_AW3D]} or pass \texttt{surface\_elevation}. This
observed-to-modeled elevation substitution requires validation for the intended
application.

The existing debris-SMB hook calls the predictor when
\texttt{processes.debris\_cover.smb.cliffs} is true. It writes the four category
fractions to \texttt{state.cliffs\_s}, \texttt{cliffs\_p}, \texttt{cliffs\_c},
and \texttt{cliffs\_u}, and their sum to \texttt{state.cliffs\_total}. The SMB
adjustment then uses
\[
 \mathrm{SMB}_{\rm out}=\mathrm{SMB}_{\rm debris}(1-f_{\rm cliff})
 +\mathrm{SMB}_{\rm clean}f_{\rm cliff}\,E_{\rm cliff},
\]
where $f_{\rm cliff}$ is the predicted total area fraction and
$E_{\rm cliff}$ is \texttt{cliff\_smb\_enhancement}. The old three-predictor
regressor can be selected by setting \texttt{model\_file} to
\texttt{cliff\_type\_models.joblib}.

\section{Configuration}
Configuration is nested under
\texttt{processes.debris\_cover.smb.cliff\_prediction} in
\texttt{user/conf/processes/debris\_cover.yaml}. Important choices are:
\begin{description}[style=nextline]
  \item[\texttt{model\_file}] Selects the serialized model bundle: \texttt{cliff\_prediction\_models.joblib} (default run015), \texttt{cliff\_prediction\_models\_run016.joblib}, \texttt{cliff\_prediction\_models\_run039.joblib} (50 m), \texttt{cliff\_prediction\_models\_run041.joblib} (100 m), or \texttt{cliff\_type\_models.joblib} (legacy three-predictor model). The bundle supplies the model grid and feature schema.
  \item[\texttt{presence\_mode}] \texttt{spatial\_stochastic} (default), \texttt{independent\_stochastic}, or \texttt{threshold}.
  \item[\texttt{probability\_calibration}] Toggle training-only Platt calibration for stochastic presence. It defaults off for run039/run041 because no calibrators were fitted for those classifiers; enabling it for those backends raises an error.
  \item[\texttt{presence\_spatial\_correlation}] Toggle spatial correlation in the presence field.
  \item[\texttt{fraction\_spatial\_correlation}] Toggle spatial correlation in magnitude fields.
  \item[\texttt{fraction\_sum\_policy}] \texttt{rescale}, \texttt{allow}, or \texttt{error} for category sums above one.
  \item[\texttt{allow\_surface\_slope\_fallback}] Allow surface descent when modeled velocity is missing or slow.
  \item[\texttt{window\_policy}] \texttt{physical} approximates 300 m/1100 m windows off-grid; \texttt{fixed\_cells} keeps 3\,$\times$\,3 and 11\,$\times$\,11.
  \item[\texttt{require\_100m\_grid}] Applies to direct array calls and requires the selected backend's model grid. The IGM state adapter resamples to that grid automatically.
  \item[\texttt{random\_seed}] Base seed for reproducible spatial fields.
\end{description}

\section{Validation and Limitations}
The implementation was smoke-tested under the igm-3.2 environment on synthetic
50, 100, and 200 m IGM grids, including the debris-SMB call path. Feature vectors
were compared to the run-015 training extractor on sampled glacier cells: the
finite-value pattern matched; remaining numerical differences were small
relative to feature scales. A full glacier run should still be checked before
scientific deployment.

The current classifier was trained with Millan2022 observed velocity but the
IGM adapter uses modeled velocity, as requested. This predictor substitution
can shift the model distribution and should be validated against mapped cliffs.
Other limitations include isotropic covariance despite flow-aligned glacier
geometry, covariance ranges estimated from a limited training sample, the
single-realization variability of stochastic output, and changes to category
marginals if the total-fraction rescaling policy is triggered.

\end{document}
